\section{The Problem: Definitions Without Structure}
\label{sec:problem}

The accepted definitions of Digital Public Infrastructure share a common defect: they describe desired outcomes without specifying the architectural preconditions required to achieve them. The standard formulation (that DPI is a set of ``shared digital systems'' built on ``open standards'' \citep{g20}) leaves open four distinct vectors for capture.

First, \textbf{Open Standards are not Open Execution.} A proprietary system can speak a standard language (like a browser displaying HTML) while remaining entirely opaque in its internal logic. The interface is public; the machine is private. (Addressed by the \textbf{CODE} test, Section~\ref{sec:tests}.)

Second, \textbf{Legal Frameworks do not guarantee Technical Rights.} A system that requires a court order to correct a database error effectively denies correction to the majority of its users. Legal remedy operates on bureaucratic time; infrastructure operates on digital time. (Addressed by the correction velocity inequality, Section~\ref{sec:dynamics}.)

Third, \textbf{Aspiration is not Compliance.} Definitions that rely on what a system ``can'' do or ``should'' be offer no resistance to authoritarian drift. Modal language creates no enforceable constraint. (Addressed by the \textbf{GOVERN} test, Section~\ref{sec:tests}.)

Fourth, \textbf{Scale is not Accountability.} A system that serves a billion users achieves universality, but universal deployment does not confer public legitimacy. Population-scale operation may simply mean population-scale capture. The ``public'' in such definitions refers to the number of subjects, not the structure of control. (Addressed by the essential services problem, Section~\ref{sec:discussion}.)

Without hard structural requirements, ``Public'' becomes a mere descriptor of scale rather than a mode of accountability. Just as classification systems become invisible as they gain acceptance \citep{bowker1999}, incorrigible infrastructure recedes into the background, becoming unchallengeable precisely as it becomes essential. This danger is acute in the transition to AI-mediated governance, where ``Infrastructure'' begins to make probabilistic judgments rather than deterministic records.